\documentclass[10pt,twocolumn]{article}

\usepackage[top=0.28in,bottom=0.28in,left=0.32in,right=0.32in]{geometry}
\usepackage[T1]{fontenc}
\usepackage[utf8]{inputenc}
\usepackage{lmodern}
\usepackage{microtype}
\usepackage{titlesec}
\usepackage{titling}

\setlength{\parindent}{0pt}
\setlength{\parskip}{0.22em}
\setlength{\columnsep}{0.18in}
\setlength{\droptitle}{-5.5em}
\pretitle{\begin{center}\large\bfseries}
\posttitle{\par\end{center}\vspace{-1.6em}}
\preauthor{}
\postauthor{}
\predate{}
\postdate{}
\titlespacing*{\section}{0pt}{0.35em}{0.08em}
\pagestyle{empty}

\title{Note for Therapy}
\author{}
\date{}

\begin{document}

\maketitle

I wrote this personal note because I am trying to understand how something real in my marriage could still exist and still not be enough to guide what happened next. I am not trying to use neuroscience to explain Camila or prove that I am right. I am trying to use the language of memory selection to understand the rupture without erasing what was real between us.

The scientific idea that stayed with me is that memory updating does not simply erase what came before. Older and newer patterns can both remain present, but only one may become selected strongly enough to guide perception and action. That distinction between presence and control is what I am trying to process.

I also wrote a formal whitepaper because I needed to separate the mechanism from the personal pain. The scientific version asks how inhibitory timing can select between competing memory patterns. The personal version asks what it means when a bond can remain present while no longer guiding a person's actions.

I may share the personal whitepaper because it says what I have been trying to say more clearly than I can say it in a call. It is also part of how I am trying to let go. Not erasing what we had does not mean staying tied to it. It means recognizing that it mattered, placing it somewhere real, and no longer waiting for it to return as the thing that organizes the present.

\vspace{0.2em}
\hrule
\vspace{0.35em}

Camila,

I was telling you yesterday that I had been reading a study about the inhibitory system, the same broad system affected by medication like yours. I wanted to understand whether changes in that system could affect which memories become active, how strongly they shape perception and how much certainty they carry in the present.

When the present activates more than one learned association, the inhibitory system has to filter competing inputs. Some input reflects what is happening now. Some comes from associative memory. Signals in the brain arrive in very tight time windows, and inhibitory signals can arrive just early enough to block one pattern while allowing another to build. Because of that, a very small modulation of timing can let one stream dominate and suppress the others, producing a winner-take-all pattern where one interpretation becomes the felt reality of the moment.

Inhibition also shapes what gets treated as signal and what gets treated as noise. It can make a single interpretation feel absolute, even if other equally salient possibilities are still present but held in the background.

There is also something called fear extinction. New experience does not simply erase an older fear structure. If an older anxious pattern keeps organizing perception, later evidence can become background instead of changing what feels true. That is what I kept feeling in our relationship. No matter how much I was present, the life we were building could still be filtered through something older.

Inhibitory circuits help organize memory, anxiety and emotional regulation. Changes in this system could affect what feels threatening and what feels freeing.

In those months after you moved into your new place, I saw us working together in a more stable way to make a home together. It felt like the present was finally being interpreted through the life we were making, and I had almost forgotten about the older patterns that kept pulling us apart.

When everything happened this week, it felt to me like something very fast and very complete took hold. A full sense of certainty appeared all at once, and it was hard for me to reconcile that with everything we had just been building. Reading this work made me think about how that kind of certainty can arise, not necessarily as a deliberate choice, but as one internal pattern becoming strong enough to organize everything in that moment.

What we recall can be pulled between more than one memory pattern, and small changes in timing or inhibition can be enough to make one pattern settle first. Once a pattern is selected, it can feel much more certain than it did a moment before.

When you said, ``it's not like I don't miss you,'' I heard that missing me was present but not selected. I heard something trying to come through, but not becoming the pattern that could guide what you were doing.

The medication structure made me think about the difference between having the medication held inside a structure and having access to it without the same structure around it. I learned that medication like this does not turn off one specific memory circuit. It broadly changes inhibition across the brain. The distinction matters because some inhibitory systems do not carry the memory itself. They help regulate which pattern is allowed to take hold.

When it was limited by weekly appointments and a small amount each week, there seemed to be more continuity in how we stayed connected to the life we were building. When that structure changed, it felt like the balance around that selection system became less stable. I am not saying the medication determines anything by itself. I am saying that changing inhibition broadly could affect timing, filtering and certainty in a system that is already carrying more than one pattern.

I see how quickly everything became final for you, but I have also seen how different states can take over and then change. I am not out to prove you wrong. I just did not want our time to end from a moment that had not had time to stabilize.

What I am trying to understand through all of this is how something real could still exist and still not be enough to hold the present in place.

For me, not erasing what we had does not mean staying connected or turning it into something that continues. It means putting it in its place and recognizing that it mattered and that it changed me.

I am not trying to hold on to it in a way that keeps us tied together. I am trying to carry it without losing myself in it.

\end{document}
